
\begin{table}[H]
    \centering
    \resizebox{\linewidth}{!}{
    \begin{tabular}{l|l|l}
        \textbf{Choice category} & \textbf{Description} & \textbf{Methods} \\
        \hline

        %===========================
        % CUT SEPARATION
        %===========================
        \multirow{5}{*}{Cut separation} 
        & \multirow{5}{*}{How to identify violated inequalities} 
        & Gomory cuts \cite{gomory_outline_1958} \\
        & & MIR / MIR+ cuts \cite{nemhauser_integer_1988} \\
        & & Flow cover cuts \cite{atamturk_cover_2005} \\
        & & Clique and lift-and-project cuts \cite{balas_mixed_1996} \\
        & & Disjunctive cuts \cite{balas_disjunctive_1979,floudas_generalized_2008} \\
        \hline

        %===========================
        % CUTTING STRATEGY
        %===========================
        \multirow{2}{*}{Cutting strategy} 
        & \multirow{2}{*}{When to apply cutting-plane rounds}
        & Root-node cutting \cite{caprara_0_1996} \\
        & & Hybrid cut+branch strategies \cite{padberg_facet_1990} \\
        \hline

        %===========================
        % SEPARATION SCOPE
        %===========================
        \multirow{4}{*}{Separation scope} 
        & \multirow{4}{*}{Where cuts are valid}
        & Global cuts \\
        & & Local (node-specific) cuts \\
        & & Path-dependent cuts \\
        & & Hybrid global/local cut handling \\
        \hline

    \end{tabular}}
    \caption{Summary of choices and strategies specific to Branch-and-Cut, structured by decision category and representative methods.}
    \label{tab:branch_and_cut_choices}
\end{table}
